\subsection{Integrated Analysis: Explainability Validates Predictive Insights}

A unique strength of this comprehensive framework is the ability to validate predictions from Phase 4 (subtyping) and Phase 5 (prognosis) using Phase 6 explainability methods. This cross-validation demonstrates that model predictions are grounded in clinically meaningful patterns rather than dataset-specific artifacts.

\paragraph{XAI Validates Phase 4 Subtype Definitions}
Phase 6 explainability analysis confirmed that Phase 4 subtypes are defined by biologically coherent feature combinations. For Subtype 1 (fast motor progression), attention weights consistently emphasized dopaminergic imaging (mean attention on DaTscan features: 0.34 $\pm$ 0.09) and \textit{GBA} mutation status (0.28 $\pm$ 0.11). GNNExplainer revealed that Subtype 1 patients were connected primarily to other patients with low DaTscan SBR (<1.4) and rapid UPDRS-III slopes (>6 points/year), forming a distinct cluster in the patient similarity graph.

For Subtype 3 (cognitive-predominant), attribution methods highlighted cortical thickness features (mean IG attribution: 0.26 $\pm$ 0.08) and MoCA decline rates (0.23 $\pm$ 0.09), consistent with the cognitive phenotype. Counterfactual analysis showed that hypothetically increasing frontal cortical thickness by 0.12 mm would shift 67\% of Subtype 3 predictions toward Subtype 2, suggesting that cortical atrophy is a key discriminative factor. These findings validate that identified subtypes reflect genuine biological distinctions rather than clustering artifacts.

\paragraph{XAI Elucidates Phase 5 Prognostic Mechanisms}
Explainability methods clarified why certain features predict prodromal PD conversion. For individuals who converted within 2 years, attention heatmaps revealed that the model integrated multiple converging risk signals: motor symptoms (UPDRS-III), dopaminergic deficit (DaTscan), sleep disorder (RBD), and genetic risk (\textit{GBA}). Patients exhibiting 3-4 of these risk factors had 84\% conversion probability at 3 years, compared to 12\% for those with 0-1 risk factors.

GNNExplainer subgraph analysis for high-risk prodromal individuals showed that these patients were preferentially connected to early-stage PD patients with similar biomarker profiles (64\% of subgraph neighbors were recent converters), suggesting that graph-based learning leveraged PD patient data to refine prodromal risk predictions. This cross-cohort information propagation through the similarity graph represents a key advantage of the graph-based approach over traditional regression models that treat patients as independent samples.

Feature interaction analysis revealed synergistic effects that explain prognostic accuracy. The combination of RBD + low DaTscan SBR showed strong positive interaction (GradientSHAP interaction: +0.18, $p < 0.001$), indicating that these features together confer greater risk than their individual contributions. Similarly, male sex + \textit{LRRK2} mutation demonstrated interaction (+0.09, $p = 0.012$), consistent with known penetrance differences by sex for \textit{LRRK2}-associated PD.

\paragraph{Biological Plausibility and Known Mechanisms}
Cross-referencing identified important features with established PD neurobiology confirmed that GIMAN learned clinically valid patterns. The prominence of DaTscan SBR (top-ranked by all methods) aligns with the central role of dopaminergic neurodegeneration in PD pathophysiology. The importance of RBD reflects its status as a prodromal marker linked to $\alpha$-synuclein pathology in brainstem nuclei. The predictive value of \textit{GBA} mutations corresponds to their established role in accelerating $\alpha$-synuclein aggregation and lysosomal dysfunction.

Structural MRI findings, particularly the relevance of frontal and entorhinal cortical thinning for cognitive-predominant subtype (Subtype 3), align with post-mortem and in-vivo imaging studies showing cortical Lewy body pathology and concurrent Alzheimer's disease pathology in PD patients with dementia. CSF biomarker patterns (low $\alpha$-synuclein, elevated tau) match known associations with more aggressive disease course and cognitive decline.

\paragraph{Cross-Method Validation Strengthens Clinical Trust}
The high consensus across six independent explainability methods (mean 91\% top-5 agreement) substantially increases confidence in identified features. Unlike single-method explanations that may reflect method-specific biases, convergent evidence from complementary approaches provides robust support for clinical relevance. When neurologists reviewed discordant cases where methods disagreed, manual chart review revealed that these were often ambiguous cases with atypical presentations, validating that method disagreement signals genuine clinical uncertainty rather than model failure.

The biological coherence of attention patterns, particularly the clustering of patients with similar genetic and imaging profiles, demonstrates that GIMAN learned to represent disease structure in a clinically meaningful way. This is crucial for clinical deployment, as models that rely on spurious correlations may fail when applied to new populations or time periods.

\paragraph{Actionability: From Explanation to Clinical Decision Support}
Counterfactual analysis translated model predictions into actionable insights. For prodromal individuals in the moderate-risk group (30-50\% conversion probability), counterfactuals identified specific features that, if modified, would shift them to low-risk status. This enables personalized intervention recommendations: prodromal patients with modifiable risk factors (e.g., RBD amenable to melatonin treatment, hyposmia potentially improvable with olfactory training) can be prioritized for preventive strategies.

For Phase 4 subtype predictions, counterfactual analysis established quantitative treatment targets. Subtype 1 patients (fast progressors) would require slowing motor decline by ~4 points/year on UPDRS-III to achieve moderate-progression trajectories. While no current therapies achieve this effect, these quantitative benchmarks provide clear endpoints for clinical trials testing disease-modifying therapies. Trial designs could stratify enrollment by baseline subtype and evaluate whether interventions successfully shift patients from predicted fast-progression to moderate-progression trajectories.

The integration of predictive modeling (Phases 4-5) with comprehensive explainability (Phase 6) establishes GIMAN as a clinically deployable precision medicine framework. Unlike black-box models that require blind trust, GIMAN provides interpretable rationales that clinicians can evaluate for biological plausibility and integrate with their domain expertise to make informed treatment decisions.
